\documentclass[11pt,reqno]{amsart}

\usepackage{amsmath,amssymb,amsthm}
\usepackage[margin=1.15in]{geometry}
\usepackage[colorlinks=true,linkcolor=blue,citecolor=blue]{hyperref}

\newtheorem{theorem}{Theorem}[section]
\newtheorem{proposition}[theorem]{Proposition}
\newtheorem{lemma}[theorem]{Lemma}
\newtheorem{corollary}[theorem]{Corollary}
\theoremstyle{definition}
\newtheorem{definition}[theorem]{Definition}
\newtheorem{question}[theorem]{Question}
\newtheorem{remark}[theorem]{Remark}

\newcommand{\spec}{\sigma}
\newcommand{\sess}{\sigma_{\mathrm{ess}}}
\newcommand{\Mm}{M}   % major radius
\newcommand{\HH}{\mathcal{H}}
\newcommand{\PP}{\mathcal{P}}
\newcommand{\BH}{\mathcal{B}(\HH)}
\newcommand{\N}{\mathbb{N}}
\newcommand{\R}{\mathbb{R}}
\newcommand{\pluf}{\pi}

\title[Ultrafilter limits in the projection lattice]
{Ultrafilter limits in the projection lattice\\ of a Hilbert space}
\thanks{The theorems of this paper are machine-checked. The verification artifact is archived at \href{https://doi.org/10.5281/zenodo.21987889}{\texttt{10.5281/zenodo.21987889}} (version~1.0); see \cite{plufLean} for the concept DOI, which resolves to the current version.}
\author{David Feldman}
\address{Department of Mathematics and Statistics, University of New Hampshire, Durham, NH}

\author{Alexander Wilce}
\address{Department of Mathematics, Susquehanna University, Selinsgrove, PA}

\date{\today}

\begin{document}

\begin{abstract}
A bounded sequence of reals has a limit along any ultrafilter on its index set. We study
a linear analogue in which the index set becomes a separable Hilbert space $\HH$, subsets
become closed subspaces, ultrafilters become maximal filters in the projection lattice
$\PP(\HH)$ (\emph{plufs}), bounded sequences become ellipsoids, and convergence becomes
the existence of arbitrarily round slices. Three answers emerge. For plufs aligned with an
atomic masa, the analogue holds, and is equivalent to the Kadison--Singer problem, resolved
by Marcus, Spielman and Srivastava; we record an elementary packaging of that theorem:
\emph{every ellipsoid in $\ell^2$ has almost-round coordinate slices along every
ultrafilter on $\N$}. Under the continuum hypothesis the analogue fails as badly as
possible: we construct a pluf $\pluf$ and an ellipsoid $E$ of eccentricity $4$ such that
every $\pluf$-slice of $E$ has eccentricity exactly $4$; indeed the set of $\pluf$-limits
of the associated quadratic form is an entire interval. Finally, on a Hilbert space of
measurable dimension the analogue holds along any normal measure, by a soft diagonal-intersection
argument developed in a companion paper. A closing section reads these results as a
fragment of ``linear model theory'': plufs are maximal consistent incidence types, the
state-space face $S_\pluf$ is the space of their completions, \L o\'s's theorem holds
exactly on the Boolean and normal-measure fragments, and the obstruction to a linear
Feferman--Vaught theory is entanglement.
\end{abstract}

\maketitle

\section{Introduction}\label{sec:intro}

Gian-Carlo Rota liked to observe that statements of discrete mathematics often admit formal
linear analogues, and that these can turn out to be false, trivial, deep, or merely
parallel; even formulating the analogue is not mechanical, despite the functorial passage
from a set to the Hilbert space it generates. We take as our discrete statement:

\begin{center}
\emph{a bounded sequence of reals converges along every ultrafilter on $\N$,}
\end{center}

\noindent and we linearize every ingredient. Throughout, $\HH$ denotes a separable
infinite-dimensional real Hilbert space (everything below works verbatim over
$\mathbb{C}$), $\PP(\HH)$ its lattice of closed subspaces, and $\BH$ the algebra of
bounded operators.

\begin{definition}\label{def:pluf}
A \emph{projection-lattice ultrafilter}, or \emph{pluf}, is a family
$\pluf\subseteq\PP(\HH)$ of nonzero closed subspaces, closed under finite intersection and
under enlargement, and maximal with these properties; equivalently, a maximal filter in
$\PP(\HH)$.
\end{definition}

Plufs exist by Zorn's lemma. The family $\pluf_L=\{M: L\subseteq M\}$ of all closed
subspaces containing a fixed line $L$ is a pluf; we call such plufs \emph{principal}.

For the linear analogue of a bounded sequence we take a positive definite bounded
symmetric form $T$ with bounded inverse, or equivalently the ellipsoid
$E_T=\{x:\langle Tx,x\rangle=1\}$; diagonal $T$ recovers bounded sequences. Write $m(E)$
and $\Mm(E)$ for the minor and major radii of an ellipsoid and $r(E)=\Mm(E)/m(E)$ for its
eccentricity. Slicing by a closed subspace never increases eccentricity.

\begin{definition}\label{def:RSP}
A pluf $\pluf$ has the \emph{round slice property} (RSP) for $T$ if
$\inf_{M\in\pluf}r(E_T\cap M)=1$.
\end{definition}

The proposed analogue is: \emph{every pluf has RSP for every $T$}. This paper establishes
the following trichotomy.

\begin{enumerate}
\item[(A)] (\S\ref{sec:pos}) Every pluf containing the coordinate subspaces of an
ultrafilter on a fixed orthonormal basis has RSP for every $T$. This is exactly the
Kadison--Singer problem \cite{KadisonSinger}, resolved by Marcus, Spielman and Srivastava
\cite{MSS}; conversely, the RSP statement for coordinate filters is equivalent to
Kadison--Singer (Remark~\ref{rem:KSpackage}).
\item[(B)] (\S\ref{sec:neg}) Assuming the continuum hypothesis, there is a pluf $\pluf$
and an ellipsoid $E$ of eccentricity $4$ such that \emph{every} $\pluf$-slice of $E$ has
minor radius $1$, major radius $4$, and eccentricity $4$. The failure is total: the set of
$\pluf$-limits of the form is the entire interval between the essential extremes
(Proposition~\ref{prop:interval}).
\item[(C)] (\S\ref{sec:meas}) If $\dim\HH$ is a measurable cardinal and $\mathcal U$ is a
normal measure on it, the coordinate filter of $\mathcal U$ is \emph{already} a pluf, and
it has RSP for every $T$, by a diagonal-intersection argument requiring no paving theorem.
Full proofs appear in the companion paper \cite{FWIII}.
\end{enumerate}

The reading in \S\ref{sec:coda} is model-theoretic: a pluf is a maximal consistent
set of \emph{incidence} conditions, the face $S_\pluf$ of states fixing it
(\S\ref{sec:slices}) is the space of completions of that type, and RSP is \L o\'s's
theorem for atomic formulas. \L o\'s survives linearization exactly on the Boolean
fragment (by Marcus--Spielman--Srivastava) and along normal measures (softly); in between,
the gap is measured by an ellipsoid.

A companion paper \cite{FWII} develops the combinatorial theory that determines
\emph{which} plufs arise from ultrafilters on bases, characterizes maximality of the
block-generated filter by selectivity of the ultrafilter, and reduces the diagonalizability
of a pluf to the existence of a single \emph{intimate} witness.
The nearest published relative of our maximal lattice filters is Question 6.39 of
Farah--Wofsey \cite{FarahWofsey}, which asks whether every maximal filter in the projection
lattice of the \emph{Calkin} algebra lifts; we discuss the bridge in
\S\ref{sec:questions}.

\section{Structure of the space of plufs}\label{sec:struct}

The facts below are elementary. They already show why the Stone--\v{C}ech analogy
breaks, before any analysis enters.

\begin{lemma}\label{lem:maximality}
A filter $\pluf$ in $\PP(\HH)$ is maximal if and only if for every $M\in\PP(\HH)$ either
$M\in\pluf$ or there is $N\in\pluf$ with $M\cap N=0$.
\end{lemma}

\begin{proof}
The filter generated by $\pluf\cup\{M\}$ consists of the subspaces containing some
$M\cap N$ with $N\in\pluf$, and is proper exactly when no such intersection vanishes.
\end{proof}

The condition $M\cap N=0$ is far weaker than $N\subseteq M^{\perp}$: two subspaces can
meet trivially while lying at angle $0$. The lattice sees incidence and forgets angle.

\begin{lemma}\label{lem:findim}
If a pluf $\pluf$ contains a finite-dimensional subspace then $\pluf=\pluf_L$ for a line
$L$. Consequently a pluf is principal if and only if $\bigcap\pluf\neq0$, in which case
$\bigcap\pluf$ is a line; and every nonprincipal pluf contains every subspace of finite
codimension.
\end{lemma}

\begin{proof}
Choose $M\in\pluf$ of least finite dimension $d$. For $N\in\pluf$, $M\cap N\in\pluf$ has
dimension at most $d$, hence exactly $d$, hence $M\subseteq N$. If $d\geq2$, any line
$L\subsetneq M$ meets every member of $\pluf$, contradicting maximality; so $d=1$ and
$\pluf=\pluf_M$. If $0\neq v\in\bigcap\pluf$, maximality puts $\R v$ into $\pluf$. For the
last claim: if $V$ has finite codimension and $V\notin\pluf$, some $M\in\pluf$ satisfies
$M\cap V=0$, forcing $\dim M<\infty$ and principality.
\end{proof}

\begin{proposition}\label{prop:noprime}
$\PP(\HH)$ has no prime filters: no proper filter $F$ satisfies
$M\vee N\in F\Rightarrow M\in F\text{ or }N\in F$. The same holds for the lattice of
closed subspaces of any real Hilbert space with an infinite orthonormal basis, and for
the subspace lattice of any real inner product space of finite dimension at least two.
\end{proposition}

\begin{proof}
Suppose $F$ is prime. Since $F$ is nonempty and upward closed, the whole space lies in
$F$. In the presence of an infinite orthonormal basis, fix three closed subspaces
$A,B,C$ that pairwise meet trivially and pairwise span: indexing the basis as $(e_n)$,
take $A$ the closed span of the even-indexed vectors, $B$ that of the odd-indexed ones,
and
\[
C=\{x:x_{2n}=x_{2n+1}\text{ for all }n\},
\]
so that $A\cap B=A\cap C=B\cap C=0$ while $A\vee B=A\vee C=B\vee C=\HH$ (for the last two,
$e_{2n+1}=(e_{2n}+e_{2n+1})-e_{2n}$ and symmetrically). Applying primeness to the three
joins puts one member of each pair $\{A,B\}$, $\{A,C\}$, $\{B,C\}$ into $F$; whichever way
the choices fall, $F$ contains two of $A,B,C$, whose intersection is $0$ --- contradicting
that $F$ is a proper filter.

A triple as above requires $\dim A=\dim B=\dim C$, so none exists in odd finite
dimension. There the following argument serves, and covers every finite dimension at
least two. A prime filter of a finite-dimensional lattice contains a finite-dimensional
member, hence is principal at a line $\R v$ by Lemma~\ref{lem:findim}. Choose
$w\notin\R v$, put $S=\operatorname{span}\{v,w\}$ and let $C$ be a complement of $S$.
Then
\[
P=\R w\oplus C,\qquad Q=\R(v+w)\oplus C
\]
satisfy $P\vee Q\supseteq\R w+\R(v+w)+C=S+C=\HH$, while $v\notin P$ and $v\notin Q$: a
relation $v=aw+c$ with $c\in C$ forces $c=v-aw\in S\cap C=0$ and hence $w\in\R v$, and
symmetrically for $Q$. Primeness applied to $P\vee Q$ would put $\R v$ inside $P$ or
inside $Q$, which it is not.

Primeness of a filter through all of $\PP(\HH)$ would furnish a dispersion-free
$\{0,1\}$-valued lattice homomorphism, forbidden for $\dim\HH\geq3$ by Kochen--Specker
\cite{KS}. The arguments above are elementary substitutes covering also $\dim\HH=2$.
\end{proof}

Topologize the set $\Pi$ of plufs by the basis $\widehat M=\{\pluf:M\in\pluf\}$; note
$\widehat M\cap\widehat N=\widehat{M\cap N}$.

\begin{proposition}\label{prop:topology}
$\Pi$ is Hausdorff and zero-dimensional; the principal plufs form a dense set of isolated
points of cardinality $2^{\aleph_0}$; and $\Pi$ is not compact.
\end{proposition}

\begin{proof}
Each $\widehat M$ is clopen, with complement $\bigcup\{\widehat N:N\cap M=0\}$ by
Lemma~\ref{lem:maximality}; two distinct plufs are separated by an $\widehat M,\widehat N$
with $M\cap N=0$. For a line $L$, $\widehat L=\{\pluf_L\}$, so principal plufs are
isolated, and they are dense since $\pluf_L\in\widehat M$ for every line $L\subseteq M$.
If $\dim M=2$ then by Lemma~\ref{lem:findim} $\widehat M$ consists exactly of the
$\pluf_L$, $L\subseteq M$: an uncountable discrete closed subspace, so $\Pi$ is not
compact.
\end{proof}

Thus $\Pi$ contains the projective space $\mathbf P(\HH)$ densely and \emph{discretely} ---
its own topology destroyed --- with no compactification structure. Since ultrafilter
limits are nothing but the compactness of $\beta\N$, the RSP cannot be had softly.

\section{Slices, gaps, and states}\label{sec:slices}

Fix a positive $T\in\BH$ with bounded inverse and set, for $0\neq M\in\PP(\HH)$,
\[
u_T(M)=\sup_{x\in M,\|x\|=1}\langle Tx,x\rangle,\qquad
\ell_T(M)=\inf_{x\in M,\|x\|=1}\langle Tx,x\rangle,
\]
so that with $T_M=P_MTP_M|_M$ we have $u_T(M)=\max\spec(T_M)$,
$\ell_T(M)=\min\spec(T_M)$.

\begin{lemma}\label{lem:radii}
$m(E_T\cap M)=u_T(M)^{-1/2}$, $\Mm(E_T\cap M)=\ell_T(M)^{-1/2}$, and
$r(E_T\cap M)=\sqrt{u_T(M)/\ell_T(M)}$.
\end{lemma}

Downward directedness of a filter gives
$\sup_{M\in\pluf}\ell_T(M)\leq\inf_{M\in\pluf}u_T(M)$, whence:

\begin{proposition}\label{prop:gap}
$\pluf$ has RSP for $T$ iff
$\sup_{M\in\pluf}\ell_T(M)=\inf_{M\in\pluf}u_T(M)$, iff for every $\epsilon>0$ some
$M\in\pluf$ has $u_T(M)-\ell_T(M)<\epsilon$.
\end{proposition}

The simultaneous statement over all $T$ is state-theoretic. Put
\[
S_\pluf=\{\varphi\text{ a state on }\BH:\varphi(P_M)=1\text{ for all }M\in\pluf\}.
\]

\begin{proposition}\label{prop:face}
$S_\pluf$ is a nonempty weak$^*$-compact face of the state space of $\BH$; it equals the
set of states with
$\sup_M\ell_T(M)\le\varphi(T)\le\inf_Mu_T(M)$ for every self-adjoint $T$; and $\pluf$ has
RSP for every $T$ iff $S_\pluf$ is a singleton, in which case its element is a pure
state.
\end{proposition}

\begin{proof}
$\overline{\lim}_\pluf T:=\inf_Mu_T(M)$ is sublinear on self-adjoints, with
$\underline{\lim}_\pluf T=-\overline{\lim}_\pluf(-T)$; Hahn--Banach supplies a linear
functional between them, which is a state, so $S_\pluf\neq\emptyset$. If $\varphi(P_M)=1$
then $\varphi(A)=\varphi(P_MAP_M)$ and
$\ell_T(M)P_M\le P_MTP_M\le u_T(M)P_M$ gives the sandwich; conversely
$\ell_{P_M}(M)=1$ forces $\varphi(P_M)=1$. A convex combination is $1$ at $P_M$ only if
both summands are, so $S_\pluf$ is a face and its extreme points are pure. The last
sentence is Proposition~\ref{prop:gap} over all $T$.
\end{proof}

The face $S_\pluf$ realizes the pluf: in the GNS representation
$(\HH_\varphi,\rho_\varphi,\xi_\varphi)$ of any $\varphi\in S_\pluf$ one has
$\rho_\varphi(P_M)\xi_\varphi=\xi_\varphi$ for all $M\in\pluf$, i.e.\ the GNS space
contains a vector lying in every subspace of the pluf. For principal $\pluf_{\R v}$ this
returns $\HH$ with $\xi=v$. We return to this reading in \S\ref{sec:coda}.

Finally, the pluf notion sits properly inside Farah and Weaver's \emph{quantum filters}
\cite{FarahWofsey,WeaverQMC}: upward-closed families of projections with
$\|p_1\cdots p_n\|=1$ for all finite subfamilies, whose maximal elements correspond
bijectively to pure states. Every filter in $\PP(\HH)$ is a quantum filter; the converse
fails because $\|pq\|=1$ sees angle where $p\wedge q\neq0$ sees incidence, and quantum
filters need not be closed under meets.

\begin{remark}[Round subspaces are abundant]\label{rem:dvoretzky}
Definition~\ref{def:RSP} is not Dvoretzky's theorem in disguise: for ellipsoids the
corresponding fact is stronger and elementary. If $Q$ has eigenvalues
$\lambda_1\ge\cdots\ge\lambda_n$, choose $c\in[\lambda_{n+1-k},\lambda_k]$ and in each
plane $\mathrm{span}(e_j,e_{n+1-j})$, $j\le k\le n/2$, a unit $f_j$ with $Q(f_j)=c$: the
$f_j$ are $Q$-orthogonal, and $Q$ restricts to $c\cdot I$ on their span --- an
\emph{exactly} spherical section of dimension $\lfloor n/2\rfloor$. The same pairing in
$\HH$ applied to $T=\mathrm{diag}(1,\tfrac1{16},1,\tfrac1{16},\dots)$ produces an
infinite-dimensional subspace on which $T$ compresses to $\tfrac12I$. What
\S\ref{sec:neg} constructs is thus a maximal filter dodging a $2^{\aleph_0}$-sized family
of perfectly round subspaces.
\end{remark}

\section{The atomic masa: Kadison--Singer as a round-slice theorem}\label{sec:pos}

Fix an orthonormal basis $(e_n)$ of $\HH$ and for $S\subseteq\N$ let
$\HH_S=\overline{\operatorname{span}}\{e_n:n\in S\}$. An ultrafilter $\mathcal U$ on $\N$
generates the filter $\{\HH_S:S\in\mathcal U\}$ (note
$\HH_S\cap\HH_{S'}=\HH_{S\cap S'}$).

\begin{theorem}\label{thm:MSS}
Let $\mathcal U$ be an ultrafilter on $\N$ and $\pluf$ any pluf containing
$\{\HH_S:S\in\mathcal U\}$. Then $\pluf$ has RSP for every $T$, and
$\lim_\pluf T=\lim_{n\to\mathcal U}\langle Te_n,e_n\rangle$.
\end{theorem}

\begin{proof}
By Proposition~\ref{prop:face}, $S_\pluf$ lies inside the set of state extensions to
$\BH$ of the pure state $\mathcal U$ on the atomic masa. By the Kadison--Singer problem
\cite{KadisonSinger}, resolved affirmatively by Marcus, Spielman and Srivastava
\cite{MSS} via Anderson's paving reformulation \cite{Anderson79}, that set is a
singleton; being nonempty, $S_\pluf$ is that singleton, namely
$\mathcal U\circ\mathcal E$ for the diagonal conditional expectation $\mathcal E$.
\end{proof}

\begin{remark}\label{rem:KSpackage}
Unwinding, RSP along coordinate filters restates paving: Anderson's theorem partitions
$\N$ into finitely many pieces on each of which $T-\mathcal E(T)$ compresses to small
norm; one piece lies in $\mathcal U$, and shrinking further makes the diagonal nearly
constant. Thus the statement
\begin{center}
\emph{for every ultrafilter $\mathcal U$ on $\N$, every ellipsoid $E$ in $\ell^2$, and
every $\epsilon>0$ there is $S\in\mathcal U$ with $r(E\cap\HH_S)<1+\epsilon$}
\end{center}
is \emph{equivalent} to the Kadison--Singer problem.
\end{remark}

Which plufs contain a coordinate filter, for which the block-generated filter is already
maximal, and whether \emph{every} pluf is so aligned, are the subject of the companion
paper \cite{FWII}: the block filter of $\mathcal U$ (blocks together with all
finite-codimension subspaces) is a pluf iff $\mathcal U$ is selective, and a pluf is
diagonalizable over a basis iff every member is \emph{intimate} with respect to it.

\section{A pluf with no round slices, under CH}\label{sec:neg}

Fix once and for all
\[
T=\mathrm{diag}\bigl(1,\tfrac1{16},1,\tfrac1{16},\dots\bigr),\qquad E=E_T,
\]
so $\spec(T)=\sess(T)=\{\tfrac1{16},1\}$, $m(E)=1$, $\Mm(E)=4$, $r(E)=4$.

\begin{definition}
A closed subspace $M\neq0$ is \emph{ample} if
$\{1,\tfrac1{16}\}\subseteq\sess(T_M)$.
\end{definition}

Since $\tfrac1{16}\le T\le1$, ampleness forces $u_T(M)=1$, $\ell_T(M)=\tfrac1{16}$, so by
Lemma~\ref{lem:radii}:

\begin{lemma}\label{lem:ample-radii}
Every ample $M$ satisfies $m(E\cap M)=1$, $\Mm(E\cap M)=4$, $r(E\cap M)=4$; in particular
ample subspaces are infinite-dimensional.
\end{lemma}

\begin{lemma}[Upward inheritance]\label{lem:upward}
Let $W\subseteq V$ be closed and $\lambda\in\{\tfrac1{16},1\}$. If
$\lambda\in\sess(T_W)$ then $\lambda\in\sess(T_V)$. Hence ampleness passes upward and its
failure passes downward.
\end{lemma}

\begin{proof}
Since $\spec(T)=\{\tfrac1{16},1\}$ we have $(T-1)(T-\tfrac1{16})=0$, whence for every
unit $x$ the two exact identities
\[
\|Tx-x\|^2=\tfrac{15}{16}\bigl(1-\langle Tx,x\rangle\bigr),\qquad
\|Tx-\tfrac1{16}x\|^2=\tfrac{15}{16}\bigl(\langle Tx,x\rangle-\tfrac1{16}\bigr).
\]
Choose a weakly null orthonormal $(x_k)$ in $W$ with $\langle Tx_k,x_k\rangle\to\lambda$.
The identity at $\lambda$ gives $\|Tx_k-\lambda x_k\|\to0$, hence
$\|T_Vx_k-\lambda x_k\|=\|P_V(Tx_k-\lambda x_k)\|\to0$: a weakly null approximate
eigensequence for $T_V$ at $\lambda$.
\end{proof}

\begin{lemma}\label{lem:finitecodim}
If $W\subseteq V$ has finite codimension in $V$ then $\sess(T_W)=\sess(T_V)$; in
particular ampleness is inherited by finite-codimension subspaces.
\end{lemma}

\begin{proof}
$P_V-P_W$ has finite rank, so $T_V$ and $T_W\oplus0$ differ by a finite-rank operator.
\end{proof}

\begin{lemma}[Escape]\label{lem:escape}
Let $V$ be ample, $N$ closed, $V\cap N$ not ample, and
$\lambda\in\{\tfrac1{16},1\}$ with $\lambda\notin\sess(T_{V\cap N})$. Then for every
$\epsilon>0$ the spectral subspace
$K^\epsilon_\lambda(V)=\mathbf1_{[\lambda-\epsilon,\lambda+\epsilon]}(T_V)\,V$ is
infinite-dimensional, $\lambda\in\sess(T_{K^\epsilon_\lambda(V)})$, and
$K^\epsilon_\lambda(V)\cap N$ has infinite codimension in $K^\epsilon_\lambda(V)$. In
particular $N$ has infinite codimension in $\HH$.
\end{lemma}

\begin{proof}
Take $\lambda=1$. A weakly null orthonormal $(x_k)$ in $V$ with
$\langle T_Vx_k,x_k\rangle\to1$ satisfies $\|T_Vx_k-x_k\|\to0$ by the estimate above,
hence $\|\mathbf1_{[1-\epsilon,1]}(T_V)x_k-x_k\|\to0$; so $K:=K^\epsilon_1(V)$ is
infinite-dimensional with $1\in\sess(T_K)$. If $K\cap N$ had finite codimension in $K$,
Lemmas~\ref{lem:finitecodim} and \ref{lem:upward} would give $1\in\sess(T_{V\cap N})$,
contrary to hypothesis.
\end{proof}

\begin{remark}\label{rem:approxeigen}
The proof uses only three properties of $K$: infinite dimension, the estimate
$\|Tx-\lambda x\|\le\delta\|x\|$ for all $x\in K$ (with $\delta$ at the caller's
disposal), and infinite codimension of $K\cap N$ in $K$. A subspace with these
properties is available without the Borel functional calculus. Since
$\lambda\in\sess(T_V)$, the identities in the proof of Lemma~\ref{lem:upward} supply an
orthonormal $(u_n)$ in $V$ with $\|Tu_n-\lambda u_n\|<\delta_n$ for any prescribed
$\delta_n>0$; taking $\sum_n\delta_n\le\delta$ and $K$ the closed span of the $u_n$,
every $x=\sum_nc_nu_n\in K$ has $|c_n|\le\|x\|$ and hence
$\|Tx-\lambda x\|\le\sum_n|c_n|\delta_n\le\delta\|x\|$.
\end{remark}

\begin{lemma}[Blocking lemma]\label{lem:blocking}
Let $G$ be a countable, downward directed family of ample subspaces and $N$ a closed
subspace such that $q\cap N$ is not ample for some $q\in G$ while $g\cap N\neq0$ for all
$g\in G$. Then there is a closed subspace $R$ with
\begin{enumerate}
\item[(i)] $R\cap N=0$;
\item[(ii)] $R\cap g$ is ample for every $g\in G$, and $R$ is ample.
\end{enumerate}
\end{lemma}

\begin{proof}
Enumerate $G=\{g_1,g_2,\dots\}$ and put $h_n=q\cap g_1\cap\dots\cap g_n$: a decreasing
sequence in $G$ (directedness), each ample, each $h_n\cap N$ not ample
(Lemma~\ref{lem:upward}). Fix $\lambda_0\in\{\tfrac1{16},1\}$ with
$\lambda_0\notin\sess(T_{q\cap N})$, hence
$\lambda_0\notin\sess(T_{h_n\cap N})$ for all $n$, and let $\lambda_1$ be the other
value. Choose $\epsilon_n\downarrow0$.

We construct unit vectors $w_n\in h_n$ recursively so that, writing
$y_n=P_{N^\perp}w_n$:
\begin{enumerate}
\item[(a)] $|\langle Tw_n,w_n\rangle-\mu_n|<\epsilon_n$, where $\mu_n=\lambda_0$ for even
$n$ and $\mu_n=\lambda_1$ for odd $n$;
\item[(b)] $\langle w_n,w_j\rangle=0$ and $\langle Tw_n,w_j\rangle=0$ for all $j<n$;
\item[(c)] $y_n\neq0$ and $\langle w_n,y_j\rangle=0$ for all $j<n$.
\end{enumerate}
All conditions in (b) and (c) other than $y_n\neq0$ are \emph{finitely many linear
constraints} on $w_n$ (self-adjointness of $T$ converts $\langle Tw_n,w_j\rangle=0$ into
$w_n\perp Tw_j$), so they carve out a finite-codimension subspace $C_n$ of any
infinite-dimensional subspace we work in.

\emph{Even $n$.} Work in $K:=K^{\epsilon_n}_{\lambda_0}(h_n)$, infinite-dimensional with
$K\cap N$ of infinite codimension in $K$ (Lemma~\ref{lem:escape}). Then
$C_n\subseteq K$ is infinite-dimensional and $C_n\cap N\subseteq K\cap N$ has infinite
codimension in $C_n$; pick a unit $w_n\in C_n\setminus N$, so $y_n\neq0$; and every unit
vector of $K$ satisfies (a) since
$\lambda_0-\epsilon_n\le T_{h_n}\le\lambda_0+\epsilon_n$ on $K$ gives
$|\langle Tw_n,w_n\rangle-\lambda_0|\le\epsilon_n$.

\emph{Odd $n$.} First pick a unit $w_n'\in K^{\epsilon_n/2}_{\lambda_1}(h_n)$ satisfying
the linear constraints of (b)--(c); then pick a unit
$z_n\in K^{\epsilon_n/2}_{\lambda_0}(h_n)\setminus N$ satisfying the same linear
constraints together with $z_n\perp w_n'$ and $z_n\perp Tw_n'$ (again finitely many
linear conditions inside an infinite-dimensional escaping subspace, by
Lemma~\ref{lem:escape}). Set $w_n=(w_n'+\eta z_n)/\|w_n'+\eta z_n\|$ for small
$\eta>0$. Conditions (b)--(c) hold for every $\eta$ by linearity; $y_n\neq0$ for all but
at most one value of $\eta$ (as $P_{N^\perp}z_n\neq0$); and
$\langle Tw_n,w_n\rangle=(\langle Tw_n',w_n'\rangle+\eta^2\langle Tz_n,z_n\rangle)
/(1+\eta^2)$ lies within $\epsilon_n$ of $\lambda_1$ for $\eta$ small. This perturbation
step replaces the orthogonal decomposition of $h_n$ along $N$, which
non-distributivity forbids.

Let $R=\overline{\operatorname{span}}\{w_n\}$ and
$R_j=\overline{\operatorname{span}}\{w_n:n\ge j\}$.

\emph{(i).} By (b) the $w_n$ are orthonormal, so each $x\in R$ has a unique expansion
$x=\sum c_nw_n$ with $(c_n)\in\ell^2$, and $P_{N^\perp}x=\sum c_ny_n$ (the series
converges since $\|y_n\|\le1$). By (c) the $y_n$ are pairwise orthogonal
($\langle y_n,y_j\rangle=\langle w_n,P_{N^\perp}y_j\rangle=\langle w_n,y_j\rangle=0$ for
$j<n$, as $y_j\in N^\perp$), and each $y_n\neq0$; so $P_{N^\perp}x=0$ forces
$c_n\|y_n\|=0$, i.e.\ $c_n=0$, for all $n$. Thus $R\cap N=0$.

\emph{(ii).} By (b), $T_R$ is \emph{exactly diagonal} in the orthonormal basis $(w_n)$ of
$R$: $\langle T_Rw_m,w_n\rangle=\langle Tw_m,w_n\rangle=0$ for $m\neq n$ (using
self-adjointness for $m<n$), with diagonal entries
$d_n=\langle Tw_n,w_n\rangle$ accumulating at both $\lambda_0$ and $\lambda_1$ by (a).
Hence $\sess(T_{R_j})\supseteq\{\tfrac1{16},1\}$ for every $j$: $R_j$ is ample. Since
$w_n\in h_n\subseteq g_j$ for $n\ge j$, $R_j\subseteq R\cap g_j$, and
Lemma~\ref{lem:upward} makes $R\cap g_j$ ample; $j=1$ gives ampleness of $R$.
\end{proof}

\begin{remark}\label{rem:exact}
Every constraint imposed in the recursion is linear, so exact orthogonality costs
nothing and no error accumulates: $T_R$ is exactly diagonal in the basis $(w_n)$. Part
(ii) consumes only the diagonal values $d_n$ and their accumulation at
$\tfrac1{16}$ and $1$.
\end{remark}

\begin{theorem}\label{thm:CH}
Assume CH. There is a nonprincipal pluf $\pluf$ every member of which is ample.
Consequently every $\pluf$-slice of $E$ has minor radius $1$, major radius $4$, and
eccentricity $4$; in particular $\pluf$ fails RSP for $T$.
\end{theorem}

\begin{proof}
Under CH enumerate $\PP(\HH)=\{p_i:i<\omega_1\}$. We build an increasing chain
$(F_i)_{i<\omega_1}$ of countable, intersection-closed families of ample subspaces with:
\begin{enumerate}
\item[($\ast$)] for each $i$, either $p_i\in F_i$ or some $g\in F_i$ has $g\cap p_i=0$.
\end{enumerate}
Set $F_0=\{\HH\}$ and take unions at limits (countable by CH). Given
$G=\bigcup_{j<i}F_j$:

\emph{Case I: $g\cap p_i=0$ for some $g\in G$.} Set $F_i=G$.

\emph{Case II: $g\cap p_i$ is ample for every $g\in G$.} Set
$F_i=G\cup\{p_i\}\cup\{g\cap p_i:g\in G\}$ (note $p_i=\HH\cap p_i$ is ample).

\emph{Case III: otherwise.} All $g\cap p_i\neq0$ but some $q\cap p_i$ is not ample.
Apply Lemma~\ref{lem:blocking} with $N=p_i$ to get $R$, and set
$F_i=G\cup\{R\}\cup\{g\cap R:g\in G\}$.

In each case $F_i$ is countable, intersection-closed, ample, and satisfies ($\ast$).
Let $\pluf$ be the upward closure of $\bigcup_iF_i$: a proper filter of ample (hence
nonzero) subspaces, every member of which is ample by Lemma~\ref{lem:upward}, maximal by
($\ast$) and Lemma~\ref{lem:maximality}, and nonprincipal since ample subspaces are
infinite-dimensional (Lemma~\ref{lem:findim}). Lemma~\ref{lem:ample-radii} gives the
radius assertions.
\end{proof}

\begin{remark}
The pluf of Theorem~\ref{thm:CH} contains no coordinate subspace $\HH_S$ with $S$ almost
contained in the evens or the odds, so its trace on the atomic masa is a filter on $\N$
that is \emph{not} an ultrafilter: the blocking in Case III is performed by subspaces
transverse to the masa. Lattice maximality never forces a pluf to decide $S$ against
$\N\setminus S$; that is why Theorem~\ref{thm:MSS} does not apply.
\end{remark}

The failure is not merely the absence of a limit; the ``limit'' is an entire interval.

\begin{proposition}\label{prop:interval}
For the pluf $\pluf$ of Theorem~\ref{thm:CH},
$\{\varphi(T):\varphi\in S_\pluf\}=[\tfrac1{16},1]$: every value between the essential
extremes is a $\pluf$-limit of $T$.
\end{proposition}

\begin{proof}
Containment $\subseteq$ is Proposition~\ref{prop:face} with $u_T\equiv1$,
$\ell_T\equiv\tfrac1{16}$ on $\pluf$. For $\supseteq$, fix $c\in[\tfrac1{16},1]$ and
define $\psi(bT)=bc$ on the line $\R T$. For $b\ge0$,
$\overline{\lim}_\pluf(bT)=b\ge bc$; for $b<0$,
$\overline{\lim}_\pluf(bT)=\tfrac b{16}\ge bc$. So $\psi$ is dominated by the
sublinear functional $\overline{\lim}_\pluf$ and extends by Hahn--Banach to a linear
$\varphi$ with $\underline{\lim}_\pluf\le\varphi\le\overline{\lim}_\pluf$ on
self-adjoints; such $\varphi$ is a state in $S_\pluf$ (Proposition~\ref{prop:face}) with
$\varphi(T)=c$.
\end{proof}

\section{Measurable dimension: a summary}\label{sec:meas}

The third act takes place on $\HH=\ell^2(\kappa)$ with $\kappa$ a measurable cardinal and
$\mathcal U$ a normal measure on $\kappa$; full proofs appear in the companion note
\cite{FWIII}, whose starting point is the following observation. Every vector of
$\ell^2(\kappa)$ has countable support, and a normal measure is closed under diagonal
intersections; together these give, for any $T\in\mathcal B(\ell^2(\kappa))$, a set
$S\in\mathcal U$ with $P_STP_S$ \emph{exactly diagonal} --- Anderson paving with error
zero, by a diagonal intersection. From this one deduces an exact dichotomy: for every
closed subspace $W$ there is $S\in\mathcal U$ with either $\HH_S\subseteq W$ or
$W\cap\HH_S=0$. Consequently:

\begin{theorem}[{see \cite{FWIII}}]\label{thm:meas}
Let $\kappa$ be measurable and $\mathcal U$ normal on $\kappa$. Then
$\Phi(\mathcal U)=\{M:\HH_S\subseteq M\text{ for some }S\in\mathcal U\}$ is a
nonprincipal, $\kappa$-complete pluf; it has RSP for every $T$, with
$\lim_{\Phi(\mathcal U)}T=\lim_{\mathcal U}\langle Te_\alpha,e_\alpha\rangle$; and
$S_{\Phi(\mathcal U)}$ is a singleton whose element is a singular, countably additive
pure state.
\end{theorem}

Three qualifications bound the scope of Theorem~\ref{thm:meas}. The collapse is local
to the normal measure:
pure states on the diagonal of $\mathcal B(\ell^2(\kappa))$ corresponding to arbitrary
ultrafilters on $\kappa$ (for instance, those concentrating on a countable set) still
require Marcus--Spielman--Srivastava for unique extension, as in Blecher--Weaver
\cite{BlecherWeaver}, Theorem 5.1. Nor is normality idle: maximality of the block
filter \emph{forces} a Q-point-type regularity on the measure (every countable-to-one
function on $\kappa$ is injective on a measure-one set), so that maximality is pinched
between this property and normality; see \cite{FWIII} for the proof and for the open
sufficiency question. The pluf space over $\ell^2(\kappa)$ remains
non-compact and prime-filter-free (the arguments of \S\ref{sec:struct} are
dimension-free). And non-diagonalizable plufs are not excluded at $\kappa$: graphs of
unitaries between complementary blocks are non-intimate there just as at $\aleph_0$
\cite{FWII}.

The state $\varphi_{\mathcal U}=\lim_{\mathcal U}\langle\,\cdot\,e_\alpha,e_\alpha\rangle$
is countably additive by $\kappa$-completeness and singular because every projection of
countable rank is annihilated; it is thus exactly the boundary object of Gleason's
theorem, which forces countably additive states on $\BH$ to be normal precisely when
$\dim\HH$ is non-measurable. The second author of \cite{BlecherWeaver} together with
Farah developed the quantum-filter side of this boundary; the lattice side, developed in
\cite{FWIII}, is that at measurable $\kappa$ --- and, below the first Ulam measurable
cardinal, \emph{only} there --- quantum-filter maximality and lattice maximality
coincide, the bridge being the excision theory of \cite{BlecherWeaver}.

\section{A model-theoretic coda}\label{sec:coda}

The results above admit a model-theoretic reading. It requires no new theorems, only a
dictionary.

Continuous model theory \cite{BYBHU,FHS} linearizes \emph{structures}: metric domains,
uniformly continuous function symbols (our operators $\HH^{\otimes n}\to\HH$),
$[0,1]$-valued predicates, $\sup/\inf$ quantifiers. Kuperberg--Weaver's quantum relations
\cite{KuperbergWeaver} linearize \emph{relations}: operator bimodules, reducing over the
atomic masa to ordinary relations on $\N$. Both keep the \emph{index} of an ultraproduct
classical. The pluf proposal linearizes the ultrafilter itself, and the present paper
computes what happens.

In the dictionary: a pluf is a maximal consistent set of incidence conditions
``$v\in M$'' --- a maximal quantifier-free partial type in the incidence fragment;
$S_\pluf$ is the space of its completions to full types over $\BH$ (pure states being
the complete types); the GNS construction is the linear ultrapower, realizing the type
(principal plufs give realized types, and the GNS vector $\xi_\varphi$ of
\S\ref{sec:slices} is the realizing element); and RSP is \L o\'s's theorem for atomic
formulas. Then:
Proposition~\ref{prop:noprime} says \L o\'s fails for disjunction --- unavoidably, by
Kochen--Specker, so any honest linear semantics is lattice-valued, as in
Takeuti--Ozawa quantum set theory \cite{Ozawa};
Proposition~\ref{prop:topology} says the compactness theorem fails for the fully
linearized logic, quantum filters being the metric repair that buys it back;
Theorem~\ref{thm:MSS} says the Boolean reduct of a type determines its completion
(\L o\'s recovered on the Boolean fragment, at the cost of \cite{MSS});
Theorem~\ref{thm:CH} says incidence-maximality is consistently far from semantic
completeness, with the gap an interval (Proposition~\ref{prop:interval}); and
Theorem~\ref{thm:meas} says \L o\'s holds outright along normal measures.

A further obstruction appears at the next step. Products of
structures require product plufs on $\HH\otimes\HH$; but the filter generated by
$\{\overline{M\otimes N}:M\in\pluf_1,N\in\pluf_2\}$ never decides an entangled subspace,
and almost all of $\PP(\HH\otimes\HH)$ is entangled. Classically $\mathcal U\times
\mathcal V$ exists canonically (Fubini); linearly, maximalization requires uncanonical
choices on entangled data. We conjecture this failure of a linear Feferman--Vaught
theorem is the structural reason no naive linear model theory exists, and note its family
resemblance to the tensor-versus-commuting distinction underlying $\mathrm{MIP}^*=
\mathrm{RE}$.

\section{Questions}\label{sec:questions}

\begin{question}[The quadrant question]\label{q:quadrant}
Diagonalizable plufs have singleton $S_\pluf$ (Theorem~\ref{thm:MSS}); under CH some pluf
has neither property (Theorem~\ref{thm:CH}). Is there a \emph{non-diagonalizable} pluf
with $S_\pluf$ a singleton? A positive answer produces a non-diagonalizable pure state
whose lattice filter is maximal --- strictly stronger than the known non-diagonalizable
states of Akemann--Weaver \cite{AW} and Koszmider \cite{Koszmider}, none of which yield
lattice filters at all (their states being singular, hence non-regular, the meet-stability
of \cite{BlecherWeaver}, Lemma 4.3 is unavailable). A negative answer makes RSP-for-all-$T$
a \emph{characterization} of diagonalizability: a lattice-level converse to
Marcus--Spielman--Srivastava.
\end{question}

\begin{question}\label{q:tstar}
CH enters Theorem~\ref{thm:CH} at exactly one point: countability of
$\bigcup_{j<i}F_j$, which makes the $h_n$ of Lemma~\ref{lem:blocking} a decreasing
$\omega$-nest. Following the template of Farah--Wofsey \cite{FarahWofsey}, Theorem 6.46
--- which derives a non-diagonalizable pure state from $\mathfrak b<\mathfrak t^*$, where
$\mathfrak t^*$ is the least length of a maximal decreasing chain in the Calkin
projection lattice --- identify the pluf-analogue of $\mathfrak t^*$ (the least length of
a maximal decreasing chain of ample subspaces, say) and derive Theorem~\ref{thm:CH} from
an inequality against it, CH becoming a corollary. Is Theorem~\ref{thm:CH} a theorem of
ZFC? Is it consistent that every pluf has RSP for every $T$?
\end{question}

\begin{question}\label{q:calkin}
Every nonprincipal pluf contains all finite-codimension subspaces
(Lemma~\ref{lem:findim}) and is invariant under finite-rank perturbations of its
projections, but not obviously under compact ones --- the gap being, once more, angle
zero. Does the image of a nonprincipal pluf generate a maximal filter in the projection
lattice of the Calkin algebra? Does every maximal filter there so arise? Either answer
bears on Question 6.39 of \cite{FarahWofsey} (whether every maximal Calkin filter lifts
to a commuting family); in particular, is the image of the ample pluf of
Theorem~\ref{thm:CH} maximal but not liftable?
\end{question}

\begin{question}\label{q:continuous}
Kadison and Singer proved non-uniqueness of pure state extensions from the continuous
masa $L^\infty[0,1]$ \cite{KadisonSinger}. Does some pluf extending the filter of
spectral subspaces of the continuous masa have non-degenerate $S_\pluf$ in ZFC? (The
countable chain $\{L^2[0,t_n]\}$ alone cannot witness this: any countable nested family
is simultaneously coordinatizable, hence intimate for the coordinatizing basis
\cite{FWII}. The full family $\{L^2[0,t]\}_{t\in(0,1]}$ generates the continuous masa
and is not coordinatizable.)
\end{question}

\subsection*{Formal verification}
This paper has been formalized and machine-checked in Lean~4 with Mathlib
(toolchain v4.28.0), as part of a verification of this series; the artifacts compile
with no unproved obligations, and every statement's axiom footprint is exactly
\texttt{propext}, \texttt{Classical.choice}, \texttt{Quot.sound}. In
Sections~\ref{sec:struct}--\ref{sec:pos}:
Lemma~\ref{lem:maximality} together with its converse (a proper, upward closed,
meet-closed family of closed subspaces satisfying the criterion is maximal);
Lemma~\ref{lem:findim} in all three clauses; Proposition~\ref{prop:noprime} in both of the
forms proved above --- for $\HH$ by the triple of subspaces, and for every space of
finite dimension at least two by the principality argument, so that no appeal to
Kochen--Specker is needed at any dimension;
Proposition~\ref{prop:topology} in all four clauses, over the topology generated by the
sets $\widehat M$; Lemma~\ref{lem:radii} in the ellipsoid form displayed, together with
the equivalence of that form of the round-slice property with its Rayleigh-quotient
form; Proposition~\ref{prop:gap}; Proposition~\ref{prop:face} in full, including
nonemptiness by Hahn--Banach, the sandwich characterization, the face property, weak\,*
compactness by Banach--Alaoglu, and the equivalence of the round-slice property for all
$T$ with singleness of $S_\pluf$; and Theorem~\ref{thm:MSS}, parametrized by a paving
hypothesis in the sense of Anderson --- the Kadison--Singer theorem of
\cite{MSS} being the quarantined classical import --- together with the elementary half
of the converse recorded in Remark~\ref{rem:KSpackage}, namely that block-witnessed
round slices for every operator yield paving.

Since Mathlib provides no state theory applicable to the bounded operators of a real
Hilbert space, states are rendered directly as continuous linear functionals that are
nonnegative on operators with nonnegative quadratic form and take the value $1$ at the
identity; all the properties used above (monotonicity, adjoint invariance,
Cauchy--Schwarz, and the compression identity at a projection of value $1$) are derived
from that definition. Section~\ref{sec:neg} has since been formalized as well, and with it every
assertion of this paper. Mathlib has no essential spectrum, no Calkin algebra, no
Fredholm theory and no Borel functional calculus; the first was defined directly by
Weyl sequences, and the last was circumvented rather than built. In place of the
spectral subspace $K^\epsilon_\lambda(V)$ of Lemma~\ref{lem:escape} the mechanized
development uses the closed span of an orthonormal sequence of approximate
eigenvectors with summable defects, which supplies exactly the three properties the
argument consumes: infinite dimension, a bound $\|Tx-\lambda x\|\le\delta\|x\|$
valid on the whole span, and infinite codimension of the intersection with $N$.
Lemma~\ref{lem:ample-radii} is verified in the ellipsoid form displayed here;
Lemmas~\ref{lem:upward}, \ref{lem:finitecodim} and \ref{lem:escape} as stated; the
blocking lemma, Lemma~\ref{lem:blocking}, in full, its recursion carried by a
dependent-choice gadget whose stage predicate refers only to the vectors already
chosen; and Theorem~\ref{thm:CH} and Proposition~\ref{prop:interval} by a
transfinite recursion in which the three cases of the construction are internalized
in a single stage lemma, so that the recursion itself never branches. The continuum
hypothesis enters as an equinumerosity hypothesis on the statements that use it and
is nowhere an axiom of the development. The mechanization also fixed two points of economy now incorporated above:
Remark~\ref{rem:exact}, that part (ii) of the blocking lemma consumes only the diagonal
values; and the reduction of Proposition~\ref{prop:interval} to the ray $\R T$. Not formalized: \S\ref{sec:coda}, and the summary
\S\ref{sec:meas}, whose contents are theorems of \cite{FWIII} and are verified
there. The verification artifact is \cite{plufLean}.

\subsection*{Acknowledgements}
We are grateful to W.~T.~Gowers for correspondence around 2000 on intimate subspaces,
which shaped the companion paper \cite{FWII}; to the late Andrew Gleason for the
suggestion, made in conversation with the first author, to consider these questions on a
Hilbert space of measurable dimension, a suggestion redeemed in \S\ref{sec:meas} and
\cite{FWIII}; and to the authors of \cite{BlecherWeaver} and \cite{FarahWofsey}, whose
published work repeatedly turned out to contain exactly the tools we needed. Portions of
these results were developed in the course of an extended dialogue with Anthropic's
Claude language model; responsibility for correctness rests with the authors.

\begin{thebibliography}{99}

\bibitem{AW} C.~Akemann and N.~Weaver, \emph{Not all pure states on $B(H)$ are
diagonalizable}, Proc. Natl. Acad. Sci. USA \textbf{105} (2008), 5313--5314.

\bibitem{Anderson79} J.~Anderson, \emph{Extensions, restrictions, and representations of
states on $C^*$-algebras}, Trans. Amer. Math. Soc. \textbf{249} (1979), 303--329.

\bibitem{BYBHU} I.~Ben Yaacov, A.~Berenstein, C.~W.~Henson and A.~Usvyatsov, \emph{Model
theory for metric structures}, in: Model Theory with Applications to Algebra and
Analysis, vol.~2, London Math. Soc. Lecture Note Ser. 350, Cambridge, 2008.

\bibitem{Bice} T.~Bice, \emph{Filters in $C^*$-algebras}, Canad. J. Math. \textbf{65}
(2013), 485--509.

\bibitem{BlecherWeaver} D.~P.~Blecher and N.~Weaver, \emph{Quantum measurable cardinals},
arXiv:1607.08505.

\bibitem{FarahWofsey} I.~Farah and E.~Wofsey, \emph{Set theory and operator algebras},
in: Appalachian Set Theory 2006--2012 (J.~Cummings and E.~Schimmerling, eds.), London
Math. Soc. Lecture Note Ser. 406, Cambridge, 2013, 63--120.

\bibitem{FHS} I.~Farah, B.~Hart and D.~Sherman, \emph{Model theory of operator algebras
I--III}, Bull. Lond. Math. Soc. / Israel J. Math. / ibid.

\bibitem{FWII} D.~Feldman and A.~Wilce, \emph{Intimate subspaces, block filters, and the
diagonalization of maximal projection filters}, in preparation.

\bibitem{plufLean} D.~V.~Feldman and A.~Wilce, \emph{pluf-lean: a Lean~4 verification of maximal filters in the projection lattice of a Hilbert space}, Zenodo, \href{https://doi.org/10.5281/zenodo.21987888}{\texttt{10.5281/zenodo.21987888}}.

\bibitem{FWIII} D.~Feldman and A.~Wilce, \emph{Projection-lattice ultrafilters at a
measurable cardinal}, in preparation.

\bibitem{KadisonSinger} R.~V.~Kadison and I.~M.~Singer, \emph{Extensions of pure states},
Amer. J. Math. \textbf{81} (1959), 383--400.

\bibitem{KS} S.~Kochen and E.~P.~Specker, \emph{The problem of hidden variables in
quantum mechanics}, J. Math. Mech. \textbf{17} (1967), 59--87.

\bibitem{Koszmider} P.~Koszmider, \emph{A non-diagonalizable pure state}, Proc. Natl.
Acad. Sci. USA \textbf{117} (2020), 33084--33089.

\bibitem{KuperbergWeaver} G.~Kuperberg and N.~Weaver, \emph{A von Neumann algebra
approach to quantum metrics / Quantum relations}, Mem. Amer. Math. Soc. \textbf{215}
(2012), no. 1010.

\bibitem{MSS} A.~Marcus, D.~A.~Spielman and N.~Srivastava, \emph{Interlacing families II:
Mixed characteristic polynomials and the Kadison--Singer problem}, Ann. of Math.
\textbf{182} (2015), 327--350.

\bibitem{Ozawa} M.~Ozawa, \emph{Transfer principle in quantum set theory}, J. Symbolic
Logic \textbf{72} (2007), 625--648.

\bibitem{WeaverQMC} N.~Weaver, \emph{Set theory and $C^*$-algebras}, Bull. Symbolic
Logic \textbf{13} (2007), 1--20.

\end{thebibliography}

\end{document}
